\section{Local kinetic energy and particle-number costs}\label{sec:energy}

The spatial estimates require two complementary comparisons. Hard exclusion
forces a kinetic cost when many particles occupy one cell or have a close
neighbor. Conversely, inserting one particle has a uniformly bounded
free-energy cost, even in an attractive external field. We establish both
comparisons, including the entropy estimates that distinguish the Gibbs
problem from its ground-state counterpart.

For a set $I$ of particle labels, write
\[
 q_I[F]=\sum_{i\in I}\int|\nabla_iF|^2,\qquad
 n_v^I=\#\{i\in I:x_i\in B_v\},\qquad
 d_i^I=\min_{j\in I\setminus\{i\}}d(x_i,x_j).
\]
The last minimum is $+\infty$ for a singleton set. Superscripts $I$ are
omitted when every label is retained. All form functions below are extended
by zero across the forbidden configurations.

\subsection{Two local consequences of the hard core}

The following arguments reproduce the local estimates of
\citet[Section~3, ``Two local kinetic inequalities'']{HS}. Their ingredients
are the two-particle coercivity used in local exclusion estimates
\citep[Proposition~10]{LundholmPortmannSolovej2015} and Neumann subdivision
\citep{LiebYngvason1998}.

\begin{lemma}[Cube energy]\label{lem:cube}
Let $Q\subset\R^3$ be a cube of side $l$, and let $F\in H^1(Q^n)$
vanish almost everywhere whenever some pair has distance at most $r$.
No boundary condition is imposed on the outer faces. Then
\begin{equation}\label{eq:cube-energy}
 \sum_{i=1}^n\int_{Q^n}|\nabla_iF|^2
 \ge c_{\rm cube}\frac r{l^3}n(n-1)\int_{Q^n}|F|^2,
 \qquad c_{\rm cube}=\frac1{64\sqrt3}.
\end{equation}
Consequently, for a torus hard-core form function and every fixed label
set $I$,
\begin{equation}\label{eq:cell-energy}
 q_I[F]\ge c_{\rm cube}r\int\sum_v n_v^I(n_v^I-1)|F|^2.
\end{equation}
\end{lemma}

\begin{proof}
Fix a point $y\in Q$. A slice $f$ vanishing on $|x-y|\le r$ satisfies
on almost every ray $x=y+s\omega$ in $Q$
\[
 |f(y+s\omega)|^2
 \le\left(\int_r^s t^{-2}\,\dd t\right)
      \int_r^s|\partial_t f(y+t\omega)|^2t^2\,\dd t
 \le r^{-1}\int_r^{R(\omega)}|\partial_t f|^2t^2\,\dd t,
\]
where convexity gives $R(\omega)\le\sqrt3l$. Rays ending before $r$
contribute zero. Integration in $s$ and $\omega$ yields
\begin{equation}\label{eq:ray-energy}
 \int_Q|f|^2\le\frac{\sqrt3l^3}{r}\int_Q|\nabla f|^2.
\end{equation}
Sobolev slicing, or approximation on the ray intervals, justifies this
calculation for $H^1$ functions.

For $n\ge2$, subdivide $Q$ into $j^3$ cubes, with
$j=\lfloor(n/2)^{1/3}\rfloor$. Then $n/16\le j^3\le n/2$.
In a sector specifying the subcube of every coordinate, at most $j^3$
coordinates occupy singleton cubes. For each of the at least $n/2$
remaining coordinates, choose a partner in its own subcube by a fixed
label ordering. Apply \eqref{eq:ray-energy} in that coordinate, with
the partner and all other coordinates fixed. Summing uses each coordinate's
kinetic energy once and gives the lower bound
\[
 \frac{rj^3}{\sqrt3l^3}\frac n2
 \ge\frac{rn^2}{32\sqrt3l^3}
\]
times the sector norm. Restriction to sectors introduces no derivative
of their indicators: this is Neumann bracketing. Sum over sectors; the
cases $n=0,1$ are immediate. This proves \eqref{eq:cube-energy} with
the stated smaller constant. Finally assign the coordinates in $I$ to
unit cells, fix the other coordinates, and apply the cube inequality
inside each occupied cell. Euclidean distance at most $r$ in a cell
lift implies torus distance at most $r$, so the required zeros persist.
Integration over the fixed coordinates proves \eqref{eq:cell-energy}.
\end{proof}

\begin{lemma}[Close-neighbor energy]\label{lem:close}
For $0<r<s<1$, every hard-core form function satisfies
\begin{equation}\label{eq:close-energy}
 q_I[F]\ge c_{\rm near}rs^{-3}
       \int\#\{i\in I:d_i^I\le2s\}|F|^2,
 \qquad c_{\rm near}=\frac{27}{262144\sqrt3}.
\end{equation}
\end{lemma}

\begin{proof}
Partition each coordinate circle into $\lfloor K/(8s)\rfloor$ equal
intervals and translate the resulting grid uniformly. Its cube side
$l$ lies in $[8s,16s]$. Two points of torus distance at most $2s$
belong to one cube with probability at least
$(1-2s/l)^3\ge27/64$. For any fixed translated grid, apply
\eqref{eq:ray-energy} to every coordinate having a partner in its
cube, choosing the partner by label order in each assignment sector.
This bounds the kinetic energy below by
$r/(\sqrt3l^3)$ times the number of such coordinates. For each
$i$ with $d_i^I\le2s$, fix a nearest partner before averaging over
grid translations. It belongs to the same cube with probability at
least $27/64$. Since $l^3\le4096s^3$, averaging proves the claim.
\end{proof}

\subsection{Entropy and admissible marginals}

We shall apply kinetic inequalities to mixed states, to their marginals,
and to states confined to a spatial region. The following estimate is the
ideal Bose entropy bound used in \citet[Section~3, ``An entropy bound'']{B}.

\begin{lemma}[Kinetic control of entropy]\label{lem:entropy}
Let $O$ be an open subset of $\Lambda_K$ with Dirichlet kinetic energy,
or the whole torus with periodic kinetic energy. For every fixed $a>0$
and every symmetric trace-one $j$-particle state $\sigma$ of finite
kinetic energy,
\begin{equation}\label{eq:entropy-bound}
 a\Tr\Bigl(\sum_{i=1}^j-\Delta_{O,i}\Bigr)\sigma-TS(\sigma)
 \ge -Tj-C_aT|O|(V^{-1}+T^{3/2}).
\end{equation}
In particular $S(\sigma)<\infty$. The statement holds for all $T>0$,
including $j=0$ with the vacuum convention.
\end{lemma}

\begin{proof}
The Gaussian image sum for the periodic heat kernel gives
$K_t(0)\le C(V^{-1}+t^{-3/2})$. Killing can only decrease its
diagonal, and hence
\[
 \Tr e^{ta\Delta_O}\le C_a|O|(V^{-1}+t^{-3/2}).
\]
The ideal Bose grand partition function at fugacity $e^{-1}$ has
logarithm
\[
 \sum_{m\ge1}\frac{e^{-m}}m\Tr e^{m\beta a\Delta_O}
 \le C_a|O|(V^{-1}+T^{3/2}).
\]
Its $j$-particle contribution is $e^{-j}$ times the canonical
partition function. The Gibbs variational inequality gives
\eqref{eq:entropy-bound}. For a general finite-energy state the same
inequality follows by finite spectral approximation, or from
nonnegativity of relative entropy to the ideal canonical Gibbs state.
This also proves entropy finiteness without an initial entropy
assumption. Hard-core states are admissible here by Dirichlet zero
extension into the free-particle space.
\end{proof}

We record a domain fact used repeatedly below. A marginal obtained by
retaining some particle coordinates of a finite-energy hard-core state
still has finite kinetic energy and satisfies the hard-core Dirichlet
condition among the retained coordinates. Indeed, first approximate a
form vector by smooth functions compactly supported in the allowed
configuration set. Slicing in the discarded variables places it in the
retained form domain with values in the $L^2$ space of the other
variables. Expansion in an orthonormal basis of that space shows that the
squared form norms of its coefficient functions sum to the integrated
form norm. Passing to the limit and then to a spectral decomposition of
the state proves the assertion for partial traces. The kinetic expectation
in the retained variables is unchanged. Symmetry of a bosonic state
places each such marginal on the corresponding symmetric subspace.

\subsection{Insertion without changing the old marginal}

The nearest-neighbor trial construction goes back to
\citet[Section~3.2]{LiebSeiringerYngvason2000}. We use the
configurationwise normalization of
\citet[Section~3, ``Normalized insertion'']{HS}, because it preserves
the old marginal exactly.

\begin{lemma}[Normalized insertion]\label{lem:insertion}
Let $w$ be a real normalized smooth one-particle orbital and $F$ an
$n$-particle hard-core form function. Suppose that on the support of
$F$ at most $k$ old positions lie within distance $2r$ of $\supp w$.
If $C\norm w_\infty^2kr^3\le1/4$, there is an admissible
$(n+1)$-particle function $\widetilde F$ of the same norm, with
the same old-coordinate position marginal, such that
\begin{equation}\label{eq:insertion-cost}
 q_{n+1}[\widetilde F]\le q_n[F]+C\bigl(
       \norm{\nabla w}_2^2+\norm w_\infty^2kr\bigr)\norm F_2^2.
\end{equation}
No symmetry in the new coordinate is asserted. Symmetry among the old
coordinates is preserved whenever $F$ has it.
\end{lemma}

\begin{proof}
Fix a nondecreasing Lipschitz cutoff $h_r$ with values in $[0,1]$,
zero up to $1.1r$, one beyond $2r$, and Lipschitz constant at most
$C/r$. Set
\[
 f_X(y)=w(y)\min_{1\le i\le n}h_r(d(y,x_i)),\qquad
 z(X)=\int f_X(y)^2\,\dd y.
\]
An empty minimum equals one. The relevant radius-$2r$ balls have
total volume at most $Ckr^3$, so $z\ge3/4$ on the support under
consideration. Almost everywhere, differentiation of the minimum
requires only a minimizing old coordinate and $y$. Thus
\begin{equation}\label{eq:insertion-gradient}
 \int\left(|\nabla_yf_X|^2+\sum_i|\nabla_if_X|^2\right)\dd y
 \le C\bigl(\norm{\nabla w}_2^2+\norm w_\infty^2kr\bigr).
\end{equation}
Here a cutoff derivative has squared size at most $C/r^2$ and is
supported in those balls; ties cause no change to the almost-everywhere
Lipschitz derivative bound.

Where $z>1/4$, put $a_X=f_X/\sqrt z$. For an old-coordinate derivative,
\[
 \partial a_X=z^{-1/2}\bigl(\partial f_X-
             a_X\langle a_X,\partial f_X\rangle_{L^2(\dd y)}\bigr).
\]
Orthogonal projection therefore bounds its squared $L^2(\dd y)$
norm by $z^{-1}\norm{\partial f_X}_2^2$. The new-coordinate
derivative has the same bound, since $z$ is independent of $y$.
Furthermore
\[
 \int a_X^2\,\dd y=1,\qquad
 \int a_X\nabla_i a_X\,\dd y=0.
\]
For $\widetilde F(X,y)=F(X)a_X(y)$ the second identity cancels
every old kinetic cross term after integrating in $y$. The first
preserves the norm and old marginal, and
\eqref{eq:insertion-gradient} proves \eqref{eq:insertion-cost}.

To define the multiplier globally replace $\sqrt z$ by
$\max(\sqrt z,1/2)$. This changes none of the integrals on the
support of $F$; a Sobolev function and its weak gradient vanish almost
everywhere on its zero set. The extra clearance $1.1r$, together with
compactly supported approximation of $F$, places the result in the
enlarged Dirichlet form domain. This proves the assertion for arbitrary
form functions, not only smooth ones.
\end{proof}

\begin{proposition}[Particle-number free-energy bounds]\label{prop:increments}
For sufficiently large $D$, uniformly in $K\ge100$,
$0<T\le D^{-1}$, and the fixed $\alpha$ window,
\begin{align}
 F_N&\ge c\alpha N,\label{eq:free-energy-lower}\\
 F_{N+k}&\ge F_N+c\alpha k \qquad(0\le k\le N),
                                      \label{eq:free-energy-increase}\\
 F_j^{(-W)}-F_{j-1}^{(-W)}&\le C\alpha
                    \qquad(1\le j\le2N).\label{eq:free-energy-insertion}
\end{align}
The last estimate holds for every bounded measurable $W\ge0$, with a
constant independent of $W$. All form domains through particle number
$2N$ are nonempty. The constants $c,C$ can be numerical; only the
large-$D$ threshold depends on the fixed window for $\alpha$.
\end{proposition}

\begin{proof}
First, all the form domains through particle number $2N$ are nonempty.
Use Lemma~\ref{lem:insertion} with $w=V^{-1/2}$, starting from the
vacuum. At every step its smallness condition is bounded by
$C(2N)r^3/V\le C\alpha^3/D^2$, which is small for large $D$.
Taking the modulus of the resulting nonzero labelled vector and
symmetrizing gives a nonzero bosonic form vector. Bounded fields do not
alter the form domain.

Reserve half the kinetic energy for Lemma~\ref{lem:entropy} and apply
Lemma~\ref{lem:cube} to the other half. Since
$\sum_vn_v(n_v-1)\ge N(D-1)$, every $N$-particle state satisfies
\[
 \mathcal F_N(\sigma)\ge
 \tfrac12c_{\rm cube}rN(D-1)-TN-CT(1+VT^{3/2}).
\]
With $T\le D^{-1}$ and $\alpha\ge1$, the negative terms are
absorbed by a fixed fraction of $\alpha N$ for large $D$. This proves
\eqref{eq:free-energy-lower}.

As in \citet[Section~3, ``Particle-number free-energy bounds'']{B},
$F_j/j$ is nondecreasing wherever the form domains are nonempty.
For a symmetric $j$-particle state let $S_i$ be the entropy of its
trace-one $i$-particle marginal, with $S_0=0$. Strong subadditivity
\citep{LiebRuskai1973Entropy} gives
$2S_i\ge S_{i-1}+S_{i+1}$; at $i=1$ ordinary subadditivity
suffices. All these entropies are finite by Lemma~\ref{lem:entropy}
and the preceding marginal-domain argument. Thus $S_i/i$ is
nonincreasing in $i$. The kinetic energy per particle is unchanged
under passage to a marginal. Applying the variational principle to
that admissible symmetric marginal proves monotonicity of $F_j/j$.
In particular
$F_{N+k}\ge(N+k)F_N/N$, which yields
\eqref{eq:free-energy-increase}.

For the upper bound, first distinguish the last label and impose symmetry
only among the first $j-1$ coordinates. Take $w=V^{-1/2}$ in
Lemma~\ref{lem:insertion}. Uniformly for $j\le2N$, the loss of norm
before normalization is at most
$Cjr^3/V\le C\alpha^3/D^2$. For large $D$, the normalized
orbital therefore satisfies
\begin{equation}\label{eq:constant-insertion}
 |a_X(y)|\le CV^{-1/2},\qquad
 \int\left(|\nabla_ya_X|^2+\sum_i|\nabla_i a_X|^2\right)\dd y
 \le C\alpha.
\end{equation}
For $p\in(2\pi/K)\Z^3$ define the isometry
\[
 v_pF(X,y)=F(X)a_X(y)e^{ip\cdot y}.
\]
The old kinetic cross terms vanish by normalization, and the new phase
adds exactly $|p|^2\norm F_2^2$, since $a_X$ is real. Old field
expectations are unchanged, while the new field is nonpositive. Hence
the form-energy increase, including the field, is at most
$(C\alpha+|p|^2)\norm F_2^2$.

These insertions form a bounded frame:
\begin{equation}\label{eq:insertion-frame}
 \sum_{|p|\le1}v_pv_p^*\le CI.
\end{equation}
Indeed, after zero extension in $y$, Parseval's identity at each $X$
gives
$\sum_p|\int a_X(y)e^{-ip\cdot y}u(X,y)\,\dd y|^2
=V\int a_X(y)^2|u(X,y)|^2\,\dd y\le C\int|u(X,y)|^2\,\dd y$.
Integrate in $X$ and restrict the frequency sum to obtain
\eqref{eq:insertion-frame}.

Let $Z_{j,\mathrm{dist}}^{(-W)}$ denote the heat trace in this
distinguished-label space, and let $E_l,\psi_l$ be an orthonormal
eigenbasis of the symmetric $(j-1)$-particle operator with field.
The scalar Jensen inequality in the spectral measure of each inserted
unit vector gives
\[
 \begin{split}
 C Z_{j,\mathrm{dist}}^{(-W)}
 &\ge\sum_{|p|\le1}\sum_l
 \langle v_p\psi_l,e^{-\beta H_{j,\mathrm{dist}}^{(-W)}}v_p\psi_l\rangle\\
 &\ge\sum_{|p|\le1}\sum_l e^{-\beta(E_l+C\alpha+|p|^2)}
 \ge cV e^{-\beta(C\alpha+1)} Z_{j-1}^{(-W)}.
 \end{split}
\]
The first inequality is the frame bound; the last uses at least $cV$
frequencies with $|p|\le1$. Positivity of the labelled heat kernel
gives
\[
 Z_j^{(-W)}\ge j^{-1}Z_{j,\mathrm{dist}}^{(-W)}:
\]
in the permutation expansion of the full symmetrizer retain the
$(j-1)!$ permutations fixing the last label, whose contributions are
nonnegative. Therefore
\[
 F_j^{(-W)}-F_{j-1}^{(-W)}
 \le C\alpha+1+T\log_+(Cj/V)\le C'\alpha,
\]
because $j/V\le2D$, $T\le D^{-1}$, and $\alpha\ge1$.
\end{proof}
